%----------
%	DATA BALANCING
%----------	

In Section~\ref{sec:eda} we performed an preliminary analysis of the dataset and find some useful and noticeable insights. For the three different types of classification tasks that we aim to tackle, we saw that classes are highly imbalanced. We could have ignored this fact and proceed with the experimentation but that wouldn't be the ideal thing to do as it will highly penalize the negative class in terms of the different metrics considered, in particular F1, precision and recall. This would make the model to learn in a deep way the positive class and neglecting the negative class, which is our main target. In order to address this problem we applied some data resampling techniques, such as down-sampling and over-sampling.

Moreover, \cite{rathpisey-2019} provide other different data resampling techniques, such as Random Oversampling (ROS),  Synthetic Minority Oversampling Technique 
(SMOTE) and Adaptive Synthetic (ADASYN). According to their work, results were improved for all oversampling techniques, while down-sampling approaches only worked for one model (naive bayes).

\begin{itemize}
    \item \textbf{Down-sampling:} The idea behind this approach is to randomly reduce the number of examples of the majority class so it equals the number of examples in the minority class.
    \item \textbf{Over-sampling:} In this case, the foundation behind this other idea is to populate the minority class with examples so it reaches the number of examples of the majority class. The different approaches that were considered are the same as presented in \cite{rathpisey-2019}. This technique is applied only to the training set after the split to avoid duplication of test examples, including also shuffling to prevent order bias.
        \begin{itemize}
            \item \textbf{Random Oversampling (ROS):} The way to proceed in this case is to generate copied examples of the minority class and randomly select them with replacement until the classes are balanced.
            
            \item \textbf{Synthetic Minority Oversampling Technique (SMOTE):} This technique was first introduced in \cite{chawla-2002}. The idea is to artificially generate new examples (synthetic examples) in the minority class, rather than simply oversampling with replacement. This method operates by first identifying the k-nearest neighbors for each sample in the minority class within the feature space, reflecting the similarities among instances. A random neighbor is then chosen, and a new synthetic sample is generated through linear interpolation between the chosen neighbor and the original sample. Specifically, the new sample $x_{new}$ is calculated using the formula:
            \begin{equation}
                x_{\text{new}} = x_{\text{original}} + \lambda \cdot (x_{\text{neighbor}} - x_{\text{original}})
            \end{equation}
            
            The python library \textit{imblearn} provides a built-in function for this task.
            
            \item \textbf{Adaptive Synthetic (ADASYN):} This idea was first presented in \cite{he-2008}. It extends the methodology of SMOTE by introducing a weighted distribution mechanism to handle imbalanced datasets more effectively. Unlike SMOTE, which generates the same number of synthetic samples for each instance in the minority class, ADASYN adjusts the number of synthetic samples based on the local density of the class. It assigns weights to the minority class samples depending on how difficult they are to classify; those that are harder to learn (i.e., more likely to be misclassified) receive higher weights. 
            
            The python library \textit{imblearn} also provides a built-in function for this task.
        \end{itemize}
\end{itemize}


